\setcounter{subsection}{0}
\section{DESARROLLO DE LA SOLUCIÓN}
\label{cap:desarrolloSolucion}

Este capítulo documenta la preparación de los datos y el modelado de CRISP-DM en las Fases~3
y~4, ligadas a los objetivos específicos tercero y cuarto: cómo se construyó el mecanismo que
convierte una lista escolar fotografiada en un kit de útiles concreto. El capítulo de diseño
(Sección~\ref{sec:impulsores} y siguientes) estableció dónde viven los componentes y por qué las
fronteras entre servicios se trazan así. Este explica cómo se modeló el comportamiento dentro de
esas fronteras, qué alternativas se descartaron y qué se reemplazó por una opción mejor durante la
construcción.

% =====================================================================
\subsection{El pipeline de resolución}
\label{sec:pipeline}

El problema se descompone en cuatro etapas secuenciales, y cada una consume la salida de la
anterior. La percepción extrae los renglones de una imagen o un PDF sin estructura. La
normalización recupera del catálogo los productos candidatos de cada renglón. La resolución decide
cuál corresponde o marca el renglón como ambiguo, y la optimización arma con los productos
resueltos los kits que pide el acudiente, que equilibran costo, fidelidad de marca, disponibilidad y personaje. La
percepción, la resolución y la optimización se alojan en el servicio de inteligencia, cuya carga
intensiva en CPU y memoria motivó su separación en el diseño. La recuperación de candidatos la
ejecuta el catálogo sobre su propia base.

La orquestación de las cuatro etapas se decidió en el ADR 0017. El pipeline se modela como trabajos
durables con estado persistido, reanudables por etapa y con una espera explícita para la revisión
humana. El registro consideró LangGraph, con persistencia de estado sobre PostgreSQL, y lo difirió:
los trabajos durables ya dan reanudación, reintentos idempotentes y auditoría con dependencias que
el equipo controla, y adoptarlo después no cambiaría el modelo de datos. Descartó también un motor
externo como Temporal o Airflow, desproporcionado para un pipeline que cabe en dos tablas, y dejó
LangGraph como camino de crecimiento si el pipeline llega a tener ramas.

El ADR 0017 preveía una biblioteca de máquinas de estados y una clave de reanudación
formada por el trabajo, la etapa y el intento. El sistema construido usa una tabla de transiciones
propia, con seis estados del trabajo, y registra cada etapa una sola vez por trabajo. El sexto,
cancelado, es terminal y se alcanza cuando el acudiente cancela su solicitud. Según el SAD,
la solución es más simple y basta para reintentar: un reintento vuelve a encolar el trabajo y
recalcula sus ítems y sus kits sin que el acudiente suba otra vez el documento ni se pague de nuevo
la percepción. Su costo es que no conserva el historial de intentos de cada etapa.

% =====================================================================
\subsection{De los prototipos al sistema construido}
\label{sec:lenguajes-dev}

La solución llegó a su forma actual por iteraciones sucesivas, cada una en su propio proyecto de
código. La prueba de concepto, del 29 de junio al 6 de septiembre de 2026, reunió por primera vez
el pipeline completo: lectura con Tesseract y PyMuPDF, raspado de tiendas, búsqueda en un índice
FAISS y un optimizador NSGA-II de tres objetivos. El 30 de julio la búsqueda pasó a PostgreSQL, con
pgvector y trigramas fusionados por rangos. Entre el 31 de julio y el 6 de agosto, un banco de
alternativas comparó almacenes vectoriales y la recuperación con y sin un modelo de lenguaje, y dejó
planteada la tarea como una correspondencia con el catálogo, sin generación de texto.

Del 29 de agosto al 21 de septiembre, los prototipos de tesis sometieron cada etapa a experimentos
con muestra y métrica: percepción, acuerdo de anotación, embeddings, bases de datos y calibración
del umbral. El 14 de septiembre, un prototipo de calibración reunió esos resultados. En paralelo, del 2 al
18 de septiembre, el producto se organizó como una flota de microservicios con librerías comunes. El
servicio de búsqueda, del 7 al 17 de septiembre, llevó la recuperación a un servicio propio, y los
frontales, del 6 al 19 de septiembre, añadieron la revisión de renglones y los perfiles de kit.

El sistema construido, desde el 22 de septiembre y hasta el 24 de octubre,
reunió ese trabajo en un monorrepositorio con
despliegue independiente por servicio (ADR 0013). Tiene seis desplegables con código, cinco
servicios de dominio y la puerta de enlace, sobre cinco bases de datos; el séptimo previsto, la
ingesta, se reemplazó por el catálogo curado desde la demanda de las listas
(Sección~\ref{sec:consolidacion}). Las cuatro etapas se implementaron en Python, incluida la
optimización con pymoo \citep{blank_pymoo_2020}, porque el cálculo de Amdahl
\citep{amdahl_validity_1967} del capítulo de diseño acotó en un 21,7\,\% como máximo la reducción
de latencia de reescribirlas en un lenguaje compilado (Sección~\ref{sec:lenguajes}). La migración
empezó por identidad, que concentraba la frontera de cumplimiento más urgente, y por catálogo, que
tenía el defecto de doble maestro (Sección~\ref{sec:monorepo}). Durante ella, las funciones de
aptitud detectaban si el código portado reintroducía una referencia directa a otro dominio
\citep{ford_building_2017}. La medición del umbral de resolución del 4 de octubre pertenece a una
línea de desarrollo que no está integrada en el sistema construido, y sus cifras se citan como
reportadas. El afinamiento de los kits, del 11 al 15 de octubre, se integró el 22 de octubre.
Entre el 19 y el 23 de octubre se añadieron la recuperación ante las caídas del catálogo y de
los procesos de fondo, las sugerencias de dirección, el alta de colegios no aliados y las colas
por RabbitMQ (ADR 0031 a 0037). El 24 de octubre se sumaron la caché y el límite de tasa con
Redis (ADR 0038).

La Tabla~\ref{tab:rag-fases} resume, por fase de CRISP-DM, el resultado de esas iteraciones para
las 84 estrategias de recuperación: las implementadas en el sistema construido, las que quedaron en
prototipos, las evaluadas y descartadas y las que no se consideraron, no se evaluaron o no aplican
al dominio. El anexo CRISP-DM, listado en la Tabla~\ref{tab:anexos}, describe cada estrategia
dentro de su fase, con su evidencia, sus resultados y su justificación. Dos rasgos dominan.
Veinte estrategias existen solo en prototipos, de modo que sus cifras no describen el sistema
construido. Y la última etapa no genera texto libre: el kit es una selección de productos del
catálogo validada por esquema, aunque la base no impone esa pertenencia con una clave foránea,
porque la línea del kit y el catálogo están en bases distintas.

\begin{table}[htbp]
\centering
\caption{Estrategias de recuperación por fase de CRISP-DM y estado en el sistema construido}
\label{tab:rag-fases}
\footnotesize
\begin{tabular}{@{}lrrrrp{6.2cm}@{}}
\toprule
\textbf{Fase} & \textbf{Impl.} & \textbf{Prot.} & \textbf{Desc.} & \textbf{Otras} & \textbf{Lo que usa el sistema construido} \\
\midrule
Comprensión del negocio & 1 & 0 & 1 & 5 & Sugerencias ancladas al catálogo, sin generación de texto \\
Comprensión de los datos & 2 & 1 & 0 & 0 & Catálogo curado desde la demanda; 175 listas de 17 colegios \\
Preparación de los datos & 13 & 6 & 1 & 4 & Metadatos y niveles del catálogo; percepción con Gemini validada por esquema; limpieza, corrección, sinónimos y atributos del renglón \\
Modelado & 11 & 9 & 5 & 10 & gemini-embedding-001; fusión aditiva del puntaje léxico y el coseno; filtros de metadatos; hasta tres opciones, una por nivel \\
Evaluación & 1 & 3 & 0 & 0 & $\tau=0{,}75$ con abstención y revisión de la lista completa \\
Despliegue & 6 & 1 & 1 & 3 & pgvector en la base del catálogo; caché de vectores en Redis; índice de conceptos en memoria; trabajos durables con espera humana \\
\midrule
Total & 34 & 20 & 8 & 22 & \\
\bottomrule
\end{tabular}
\par\smallskip
{\footnotesize Impl.: implementada; Prot.: solo en prototipos; Desc.: evaluada y descartada; Otras: considerada sin evaluar, no considerada o no aplica al dominio.}
\end{table}


% =====================================================================
\subsection{Percepción documental}
\label{sec:percepcion-dev}

La primera etapa recibe una fotografía o un PDF sin estructura garantizada y debe producir los
renglones con su cantidad y su descripción. El 1 de septiembre de 2026, los prototipos de tesis
compararon cinco enfoques: un flujo clásico de reconocimiento óptico de caracteres (PyMuPDF,
Tesseract y un analizador de expresiones regulares), Google Vision, GOT-OCR2.0, Docling y el modelo
multimodal Gemini. Se evaluaron con F1 sobre un conjunto de oro propio, en la línea de los bancos de
OCR para modelos multimodales \citep{liu_ocrbench_2024} y de extracción clave-valor
\citep{shen_ocr_2026}. La razón para preferir un modelo multimodal es estructural. El
reconocimiento con reglas depende de que el documento tenga una disposición regular de renglones, y
buena parte del corpus del objetivo primero son fotografías de listas manuscritas o impresas,
tomadas con el teléfono del acudiente y sin luz controlada. Un modelo que generaliza sobre la
disposición visual, en vez de suponer una, era la única familia compatible con esa variabilidad.
Las cifras están en la Sección~\ref{sec:percepcion}; la elección no tiene un ADR propio.

En el sistema construido, cada página se convierte en imagen y se envía a Gemini con un prompt en
español y el esquema de la respuesta, con salida JSON y temperatura 0. Si la respuesta no es JSON o
no cumple el esquema, el trabajo falla con su motivo: no hay reintento escondido ni renglón
inventado. Solo la falla del transporte tiene un reintento, declarado y único.

El nombre del menor es un dato
protegido por la Ley 1581 de 2012 \citep{congreso_ley1581_2012}, y la propuesta inicial era
anonimizarlo antes de guardar la lista y antes de invocar el modelo externo; los prototipos lo
intentaron con expresiones regulares y con un reconocedor de entidades. El sistema construido lo
protege de otra forma, por minimización: no lo recoge. Según PV-03 del SRS, no pide el nombre legal del menor, el
apodo con que el acudiente lo identifica es opcional y se cifra en reposo, y la percepción extrae
solo los ítems y el grado. Es mejor porque un dato que no se recoge no puede llegar a los registros
de operación, a los conjuntos de evaluación ni al modelo, y no depende de que un reconocedor lo
detecte. Queda un caso: si el nombre figura en el documento cargado, viaja con la página hacia el
proveedor, y si el modelo lo lee como renglón, queda 30 días en la lectura que guarda Redis. Su
depuración, que pide PE-06, se delimitó como parte del endurecimiento productivo previo al
despliegue con familias (Secciones~\ref{sec:defectos} y~\ref{sec:trabajo-futuro}), porque la
evaluación del trabajo usó listas con estudiantes ficticios.

% =====================================================================
\subsection{Normalización y resolución}
\label{sec:normalizacion}

Cada renglón extraído se busca en el catálogo con recuperación híbrida: una señal léxica, que
captura coincidencias exactas de marca y referencia, y una señal densa, que captura equivalencias de
significado entre descripciones redactadas de forma distinta
\citep{reimers_sentence-bert_2019,karpukhin_dense_2020}. Del marco de generación aumentada por
recuperación \citep{lewis_retrieval-augmented_2020}, esta etapa toma solo el componente de
recuperación. Antes de buscar, el texto se limpia, se
corrige la ortografía con el léxico del catálogo, se reduce al núcleo en singular y se separan los
atributos impuestos por la lista, como el color, el tamaño o la marca; los sinónimos viven en cada
concepto del catálogo.

La
prueba de concepto guardó los vectores en un índice FAISS \citep{johnson_billion-scale_2021} hasta
el 30 de julio. El sistema construido guarda el vector de 1024 dimensiones de cada producto en el
PostgreSQL del catálogo, con pgvector (Sección~\ref{sec:dominio}), junto con el modelo y el texto
del que se calculó. Así el vector se respalda con el producto, sin un segundo almacén
que sincronizar \citep{kleppmann_designing_2017}, y una huella del texto indica cuándo recalcularlo.
Los archivos versionados del catálogo no lo llevan: un proceso por lotes lo calcula contra la base.
La búsqueda se hace en memoria, con el coseno calculado en el proceso: los índices HNSW y de
trigramas existen, pero no se consultan. El SAD advierte que harán falta antes de llegar a los
50.000 productos del escenario ESC-E4.

La comparación con MTEB
\citep{muennighoff_mteb_2023} eligió e5-large, de 560 millones de parámetros, que habría que alojar
y servir. El sistema construido usa gemini-embedding-001, del mismo proveedor que la percepción. La
ventaja es operativa (no hay modelo que alojar ni GPU) y no una mejora de calidad demostrada, porque
los dos no se compararon sobre el corpus. El vector del renglón se calcula en cada solicitud y el de
cada concepto es el promedio de sus productos. Si el servicio de embeddings falla, la búsqueda sigue
solo con la señal léxica y el resultado lo advierte.

Los prototipos
fusionaban las dos señales por rangos (RRF) y reordenaban con un codificador cruzado. Sobre los
2.710 pares del conjunto de calibración, el área bajo la curva fue de 0,728 para la señal léxica,
0,805 para la densa, 0,739 para RRF y 0,880 para el reordenador. El sistema construido resuelve en
una sola fase: un puntaje léxico anclado en el núcleo del renglón más 0,15 veces el coseno, con un
rescate por coseno de 0,86 o más cuando el puntaje léxico no supera 0,55. A cambio, el
administrador de Tukit revisa completa cada lista nueva antes de publicarla (PE-04), en cuatro horas
o menos en temporada (ML-09). El SDD da tres razones. No depende de un reordenador ni de una GPU;
corrige el error que hacía tomar «chaquiras» por otro producto; y en la corrida del servicio de
búsqueda anterior, 1.407 de 2.176 ítems (64,7\,\%) necesitaban de todos modos a una persona. La
revisión se hace una vez por colegio y grado y se reutiliza (NR-07). Las constantes de la fusión (0,15; 0,86; 0,55) proceden de ajustes sobre casos de error
concretos, como el de las «chaquiras», y no de un barrido; la Tabla~\ref{tab:constantes} registra el
origen y el estado de cada una.

La estructura del catálogo hace el resto. Sus 149 conceptos, 145 de ellos activos, agrupan los
productos equivalentes; en el levantamiento de listas, 25 conceptos cubren el 80,97\,\% de los
renglones con concepto. Los atributos y los niveles distinguen
las variantes, y el sistema se abstiene cuando el producto contradice un atributo impuesto. Cada
atributo se compara a su manera (exacto, como mínimo, con tolerancia o por equivalencias), y el
que el producto no declara no cumple (ADR 0034). Un grafo
de sustitución entre productos se evaluó y se descartó: subía la cobertura del 41,3\,\% al
65,3\,\%, pero bajaba la precisión del 86,7\,\% al 83,5\,\%. La tabla de sustitución dirigida entre
conceptos curados no tiene datos y no se consulta (Sección~\ref{sec:consolidacion}).

\subsubsection{Umbral de resolución}

El umbral pasó por tres valores. El servicio de búsqueda anterior fijó $\tau=0{,}85$ y un margen
$\delta=0{,}05$ con siete ejemplos sintéticos. El prototipo de calibración lo situó en $\tau=0{,}630$
sobre el puntaje del reordenador, con 542 ítems, y así quedó en el ADR 0009. El sistema construido
usa $\tau=0{,}75$ sobre su puntaje propio, sin margen, como valor inicial de un ciclo de calibración
con revisión humana (Sección~\ref{sec:ciclos-calibracion}). Su efecto está contenido porque el umbral
solo señala renglones dentro de una lista que de todos modos se revisa completa, de modo que un
$\tau$ mal ajustado cuesta tiempo de revisión y no un producto equivocado en el carrito. El valor
vigente y su evidencia se registran en la Tabla~\ref{tab:constantes}\pendienteEvidencia{valor final de tau sobre el puntaje del sistema construido}.

\subsubsection{Algoritmo de resolución de un renglón}
\label{sec:algoritmo-resolucion}

La resolución es el algoritmo propio de OE3, y se presenta con la estructura del curso de
algoritmos del departamento: análisis del problema, diseño de entradas y salidas, estrategia,
seudocódigo y complejidad con comparación frente al método exhaustivo.

\paragraph{Análisis del problema.} Dado un renglón de la lista, por ejemplo «pompones de colores
x 20», hay que decidir a cuál de los conceptos canónicos del catálogo corresponde, es decir, a qué
clase de producto equivalente se refiere, o declarar que no corresponde a ninguno. Dos errores
motivaron el diseño. El primero es el de las «chaquiras de colores», que una regla de contención
literal asignaba al concepto «colores» (lápices de colores) porque la palabra aparecía en ambos. El
segundo es el de los sinónimos que ninguna lista contiene, como «estuche» por «cartuchera». La
solución tenía que corregir el primero sin perder el segundo y sin que la señal semántica pudiera
empeorar un acierto léxico.

\paragraph{Diseño: entradas y salidas.} La Tabla~\ref{tab:es-resolucion} define los tipos.

\begin{table}[H]
\centering\footnotesize
\caption{Entradas y salidas del algoritmo de resolución de un renglón.}
\label{tab:es-resolucion}
\begin{tabularx}{\textwidth}{@{}p{2.6cm}p{3.6cm}X@{}}
\toprule
Elemento & Tipo & Descripción \\
\midrule
Entrada $r$ & cadena de texto & Renglón de la lista, con cantidad y empaque \\
Entrada $v_r$ & vector real de 1.024 componentes, opcional & Embedding del renglón; si falta, el algoritmo usa solo la señal léxica \\
Entrada $\mathcal{C}$ & conjunto de conceptos, cada uno con nombre, lista de sinónimos y vector $v_c$ & 149 conceptos, 145 activos \\
Parámetros & tres números reales en $[0,1]$ & Aporte del coseno $\gamma=0{,}15$, umbral de rescate $\theta=0{,}86$ y piso del rescate $\phi=0{,}55$ \\
Salida & terna (concepto o nulo, nombre, puntaje en $[0,1]$) & Mejor concepto y su puntaje; un puntaje menor que $\tau$ remite el renglón a revisión \\
\bottomrule
\end{tabularx}
\end{table}

\paragraph{Estrategia algorítmica.} El método es una búsqueda del máximo puntaje sobre los
conceptos, con un puntaje en tres ramas que se evalúan en orden. Primero se extrae el núcleo del
renglón, que es la primera palabra útil después de quitar cantidades y empaques, porque en español
el núcleo del sintagma va primero. Si el núcleo está en la cabeza del nombre o de un sinónimo del
concepto, el puntaje es $0{,}6+0{,}4\,J+\gamma\cos$, donde $J$ es el parecido de Jaccard entre las
palabras significativas y $\cos$ es el coseno entre $v_r$ y $v_c$, recortado en 0. Si el núcleo
aparece en otra posición, el puntaje es $J+\gamma\cos$. Si no aparece y el coseno alcanza $\theta$, el
rescate asigna $\phi+0{,}4(\cos-\theta)/(1-\theta)$, que parte del piso $\phi$, apenas sobre el corte de revisión de
0,5. En cualquier otro caso el puntaje es 0. Como $\gamma\cos\ge 0$, el puntaje híbrido nunca es
menor que el léxico, que es la garantía que la tesis afirma, y el umbral alto de rescate impide que
un parecido de tema, como el de «chaquiras» con «colores», se cuele.

\begin{algorithm}[H]
\caption{Resolución de un renglón contra el catálogo de conceptos}
\label{alg:resolucion}
\KwEntrada{renglón $r$, vector $v_r$ (opcional), conceptos $\mathcal{C}$, parámetros $\gamma,\theta,\phi$}
\KwSalida{concepto $c^*$, puntaje $p^*$}
$n \leftarrow$ núcleo$(r)$; $(c^*, p^*) \leftarrow (\text{nulo}, 0)$\;
\ForEach{concepto $c \in \mathcal{C}$}{
  $\cos \leftarrow \max(0, \text{coseno}(v_r, v_c))$ si hay $v_r$ y $v_c$; en otro caso $0$\;
  \ForEach{candidato $k \in \{\text{nombre}(c)\} \cup \text{sinónimos}(c)$}{
    \uIf{$n$ está en la cabeza de $k$}{ $p \leftarrow \min(1,\ 0{,}6 + 0{,}4\,J(r,k) + \gamma\cos)$ }
    \uElseIf{$n$ está en alguna palabra de $k$}{ $p \leftarrow \min(1,\ J(r,k) + \gamma\cos)$ }
    \uElseIf{$\cos \ge \theta$}{ $p \leftarrow \phi + 0{,}4\,(\cos-\theta)/(1-\theta)$ }
    \Else{ $p \leftarrow 0$ }
    \lIf{$p > p^*$}{ $(c^*, p^*) \leftarrow (c, p)$ }
  }
}
\Return{$(c^*, p^*)$}\;
\end{algorithm}

\paragraph{Complejidad y comparación con el método exhaustivo.} Sea $|\mathcal{C}|$ el número de
conceptos, $s$ el número medio de candidatos por concepto (nombre y sinónimos), $w$ el número de
palabras de un candidato y $d=1.024$ la dimensión del vector. El coseno se calcula una vez por
concepto, con costo $O(|\mathcal{C}|\,d)$; la comparación léxica cuesta $O(w^2)$ por candidato con
conjuntos pequeños. El total por renglón es $O(|\mathcal{C}|\,(d+s\,w^2))$, lineal en el
tamaño del catálogo de conceptos. Con 149 conceptos y $d=1.024$, el coseno requiere unas 152.000
multiplicaciones por renglón (149 por 1.024). El método exhaustivo de referencia compara el
renglón con cada producto del catálogo en lugar de cada concepto; con 1.030 productos con precio
(Sección~\ref{sec:recuperacion}), su costo es $O(N\,d)$ con $N=1.030$, unas 1,05 millones de
multiplicaciones, es decir, 6,9 veces más que la resolución por conceptos. Esa relación es el cociente
$N/|\mathcal{C}|$ entre productos y conceptos, y crece con el catálogo, mientras que el número de
conceptos crece más despacio porque agrupa productos equivalentes. Ambos métodos son lineales, de
modo que la ventaja no es asintótica sino de constante. La razón principal para agrupar por conceptos
es la abstención: el exhaustivo por producto no distingue entre un producto ausente y un
concepto que el catálogo no cubre.

% =====================================================================
\subsection{Optimización del kit}
\label{sec:optimizacion-dev}

Con los productos resueltos, el kit se arma sobre las líneas comprables con cuatro objetivos: costo,
fidelidad de marca, disponibilidad y personaje. Ninguna ponderación fija representa a todos los
acudientes, así que el sistema construye el frente de Pareto en vez de optimizar una función de
utilidad única \citep{zitzler_multiobjective_1999}.

La prueba de concepto usaba NSGA-II
\citep{deb_fast_2002} con una población de 60, 100 generaciones, un tope de 50 ítems y 30 segundos,
y fundía la calidad con el personaje. En el banco de alternativas, sobre nueve instancias
sintéticas, quedó a un 0,07\,\% del óptimo, frente al 11,92\,\% de un algoritmo voraz. El sistema
construido conserva el algoritmo, con pymoo \citep{blank_pymoo_2020}, y cambia tres cosas. Calcula
el frente exacto cuando hay 4.096 combinaciones o menos; separa el personaje como objetivo propio; y
hace crecer la población y las generaciones con la lista, con una semilla fija. El servicio no
arranca sin pymoo.

Del frente salen siempre el Ahorrador y el Premium. El Fan sale solo si el acudiente pidió un
personaje, y «A tu Medida», el de mejor marca que cabe en el tope o, sin tope, a mitad de camino
entre los otros dos, solo si eligió el punto medio de gasto o declaró un tope (ADR 0021). Ofrecer
pocas opciones bien diferenciadas facilita decidir
\citep{thaler_nudge_2008,iyengar_when_2000,schwartz_paradox_2003} (NR-02), y un kit que nadie
pidió es ruido para el acudiente y trabajo para el revisor. La regla de recomendación marca el Fan
cuando hay personaje y «A tu Medida» cuando se armó; si no, el Premium cuando el acudiente pone la
calidad primero y acepta un sobreprecio del 10\,\%, y en los demás casos el Ahorrador. La regla sale
de la encuesta del Dossier OE1, en la que el 70\,\% de los acudientes no pagaría más de un 5\,\% de
sobreprecio. El tope es el presupuesto por hijo, sin ampliarlo por el sobreprecio aceptado. Si el
kit marcado no cabe, se marca el de mejor marca que sí cabe y, solo si ninguno cabe, el más barato,
con la razón y su costo (ADR 0020). El presupuesto acota la recomendación, no el frente: los demás
kits siguen a la vista.

Los prototipos inferían las
preferencias con reglas y, de forma opcional, con un modelo de lenguaje local. El sistema construido
las pide como opciones cerradas por hijo: cuánto gastar, el presupuesto, el sobreprecio aceptado y
el personaje (PS-01, PS-02 y PS-04). El SDD lo justifica en tres razones: la preferencia queda
explícita y auditable, no depende de un modelo y el acudiente puede cambiarla. Los pesos por omisión
son 0,35, 0,30, 0,25 y 0,10, y salen de la pregunta P21 de la encuesta normalizada (42 de 60
encuestados pagan como máximo un 5\,\% más y el personaje quedó último en la prioridad). Los de OP-04 del
SRS v6.1, 0,285, 0,273, 0,255 y 0,187, se contrastan con los anteriores en la calibración de la
utilidad multiatributo (Sección~\ref{sec:optimizacion-cal})\pendienteEvidencia{pesos MAUT finales}. El
hipervolumen del frente, indicador de calidad de un conjunto de soluciones no dominadas
\citep{zitzler_multiobjective_1999}, se calcula sobre listas reales en la
Sección~\ref{sec:optimizacion}\pendienteEvidencia{hipervolumen sobre listas reales}.

% =====================================================================
\subsection{Constantes numéricas del sistema}
\label{sec:constantes}

En esta sección se registra cada constante numérica que condiciona el comportamiento del sistema,
para que ninguna quede sin justificación. Para cada una, la Tabla~\ref{tab:constantes} indica su
valor, el módulo donde vive, su origen y su estado. El origen es de uno de tres tipos: una
referencia de la literatura, un requisito del SRS con su razón, o un experimento de los
prototipos o de la calibración. El estado es «vigente» cuando el valor tiene una justificación
registrada y «en ciclo de calibración» cuando el valor actual es el de partida de un ciclo de
revisión humana (Sección~\ref{sec:ciclos-calibracion}) cuyo resultado final se reporta en la
Sección~\ref{sec:optimizacion-cal}. Los valores proceden del código del repositorio y de la tabla de
parámetros del documento de algoritmos del proyecto.

{\footnotesize
\setlength{\tabcolsep}{3pt}
\renewcommand{\arraystretch}{1.1}
\begin{longtable}{@{}p{3.1cm}p{2.3cm}p{2.5cm}p{4.7cm}p{1.6cm}@{}}
\caption{Constantes numéricas del sistema, con su origen y su estado.}\label{tab:constantes}\\
\toprule
Constante & Valor & Módulo & Origen & Estado \\
\midrule
\endfirsthead
\toprule
Constante & Valor & Módulo & Origen & Estado \\
\midrule
\endhead
\bottomrule
\endfoot
\multicolumn{5}{@{}l}{\emph{Percepción}}\\
Tope de páginas & 40 & percepción & Experimento: el PDF más largo del corpus tiene 39 páginas & Vigente \\
Temperatura del modelo & 0 & percepción & Decodificación voraz, para reducir la variabilidad entre lecturas & Vigente \\
Tiempo límite por llamada & 40 s & percepción & Una página típica tarda unos 5 s; 40 s indica una llamada colgada & Vigente \\
Reintentos y espera & 1 reintento, 3 s & percepción & Solo ante fallas de transporte, para no ocultar respuestas mal formadas & Vigente \\
\midrule
\multicolumn{5}{@{}l}{\emph{Normalización y resolución}}\\
Dimensión del vector & 1.024 & catálogo & Fija la columna de la base; se pide explícitamente al modelo de \emph{embeddings} & Vigente \\
Lote de \emph{embeddings} & 100 textos & catálogo & Una petición por lote, para evitar cuerpos de petición excesivos & Vigente \\
Umbral de reconocimiento & 0,5 & catálogo & Corte por debajo del cual el renglón no corresponde a ningún concepto; fijado al corregir el caso «chaquiras» & Vigente \\
Aporte del coseno $\gamma$ & 0,15 & catálogo & Ajuste sobre casos de error; pequeño para que el coseno desempate y no decida & En ciclo de calibración \\
Umbral de rescate $\theta$ & 0,86 & catálogo & Ajuste sobre casos de error; separa un sinónimo de un parecido de tema & En ciclo de calibración \\
Piso del rescate $\phi$ & 0,55 & catálogo & Apenas sobre el corte de 0,5, para que un rescate entre a revisión & En ciclo de calibración \\
Umbral de autoaceptación $\tau$ & 0,75 & inteligencia & Requisito PE-04 del SRS; valor inicial anterior a la calibración (0,85 y 0,630 en versiones previas); ver Tabla~\ref{tab:cal-optim} & En ciclo de calibración \\
Tolerancia de medidas & 0,25 y rango 0,85 a 1,35 & catálogo & Ajuste con las auditorías de los kits (Tabla~\ref{tab:kits}) & Vigente \\
Rango de precio en tiendas & \$1.000 a \$400.000 & catálogo & Descarta accesorios sueltos y lotes institucionales & Vigente \\
\midrule
\multicolumn{5}{@{}l}{\emph{Optimización}}\\
Límite de enumeración & 4.096 ($2^{12}$) & inteligencia & Por debajo de ese tamaño del espacio el frente exacto se calcula en milisegundos; por encima se usa NSGA-II & Vigente \\
Población & $\min(200,\max(40,12n))$, par & inteligencia & Escala con las $n$ líneas decidibles; el par evita desperdiciar un individuo & Vigente \\
Generaciones & $\min(120,\max(30,20\log_{10}|\mathcal{E}|))$ & inteligencia & Escala con el logaritmo del espacio $|\mathcal{E}|$ de búsqueda & Vigente \\
Cruce SBX & $p_c=0{,}9$, $\eta_c=15$ & inteligencia & Valores estándar de \citet{deb_simulated_1995} & Vigente \\
Mutación polinomial & $p=1$, $\eta_m=20$ & inteligencia & Valores estándar de \citet{deb_combined_1996} & Vigente \\
Semilla & 0 & inteligencia & Reproducibilidad del frente entre ejecuciones & Vigente \\
Pesos por omisión (costo, marca, calidad, personaje) & 0,35; 0,30; 0,25; 0,10 & inteligencia & Encuesta, pregunta P21 (n=60), normalizada; ver Tabla~\ref{tab:cal-optim} & En ciclo de calibración \\
Sobreprecio para Premium & 10\,\% & inteligencia & Encuesta: 42 de 60 pagan como máximo un 5\,\% más; el 10\,\% marca a quien acepta pagar más & Vigente \\
Opciones por renglón & 3 & catálogo & Pocas opciones diferenciadas facilitan la decisión \citep{iyengar_when_2000} & Vigente \\
\midrule
\multicolumn{5}{@{}l}{\emph{Requisitos de desempeño y operación}}\\
Tiempo de generación del kit & 8 s & SRS (DP-03) & Requisito de desempeño; se contrasta en la Sección~\ref{sec:optimizacion} & Vigente \\
Ventana de revisión & 4 h en temporada & SRS (ML-09) & Requisito: el administrador revisa cada lista nueva antes de publicarla & Vigente \\
Disyuntor & 5 fallas, 30 s & biblioteca común & Patrón de resiliencia; 30 s es el escenario ESC-R2 del ATAM & Vigente \\
Caché de lecturas & 30 días & inteligencia & Evita pagar de nuevo la percepción de un mismo documento & Vigente \\
Caché de vectores & 24 h & catálogo & Evita recalcular el vector de un renglón repetido & Vigente \\
\end{longtable}
}

% =====================================================================
\subsection{Aplicación web y pago}
\label{sec:web}

El acudiente usa una aplicación web progresiva. Siguiendo las pautas de interacción con sistemas de
IA \citep{amershi_guidelines_2019}, la interfaz muestra el estado de cada renglón en nueve clases y
dice «Lo está revisando una persona» cuando la lista espera la revisión. El acudiente elige uno de
sus kits y puede cambiar el producto de una línea por el sugerido o por una de sus alternativas:
el cambio se guarda en su solicitud, se aplica al leer el kit y se cobra (ADR 0019). La dirección
de entrega se completa con las sugerencias de Google Places y, sin ese servicio, se escribe a
mano (ADR 0035).

El pago se abre en el widget de Wompi, una ventana sobre la página del kit o de «Mis pedidos»
(ADR 0024). Con tarjeta o Nequi el acudiente no sale de Tukit; con PSE, el banco toma la pestaña y
Wompi lo devuelve a la página de resultado. Comercio firma los datos del widget y guarda solo la
referencia, el identificador, el estado y el medio de la transacción (ADR 0010). La confirmación
llega por el aviso de la pasarela y por la consulta de la transacción cuando el acudiente regresa,
y las dos pasan por el mismo procedimiento idempotente (ADR 0025). El acudiente elige el medio dentro del widget, de modo que Tukit no da a PSE la prioridad que
preveía el ADR 0010. El flujo de pago se verificó en el entorno de pruebas (\emph{sandbox}) de la
pasarela con tarjeta, Nequi y PSE\pendienteEvidencia{pago de extremo a extremo en el sandbox de Wompi}. El cobro con
dinero real depende del despliegue productivo y se delimitó como trabajo futuro
(Sección~\ref{sec:trabajo-futuro}).

% Contenido que llega de otras sesiones (calibracion y MAUT; NL2C).
\subsection{Calibración de la optimización multiobjetivo}
\label{sec:optimizacion-cal}

En esta sección se describen la calibración del optimizador y el método de utilidad multiatributo que ordena los kits del frente de Pareto.

\subsection{Traducción de preferencias en lenguaje natural a restricciones}
\label{sec:nl2c}

En esta sección se describe cómo una preferencia escrita en lenguaje natural se traduce a restricciones estructuradas para el optimizador.

